\subsection{Hubbard–Stratonovich transformation order by order}
\label{p:sec:hs}

A Hubbard--Stratonovich transformation~\cite{Stratonovich:1957:MethodCalculatingQuantum,Hubbard:1959:CalculationPartitionFunctions} trades an interaction that is the square of some composite, such as $\lambda f(\phi)^2$, for a Gaussian integral over an auxiliary field coupled linearly to that composite.
For example, the quartic vertex $\lambda \phi^4$ becomes a cubic coupling $\sigma \phi^2$ plus a quadratic term for the auxiliary field, and integrating the auxiliary field out gives back the original interaction exactly (see \cref{p:app:scalar-hs}).
In closed string field theory the interaction at order $\gs^{n-2}$ is the vertex $\vertex{n}$ evaluated on $n$ copies of the string field.
The natural composites are thus the tensor powers $\Psi^k := \Psi^{\otimes k}$, and the auxiliary fields replacing them are $k$-string fields $\Sigma_k \in \Hn{k}$.
This section carries out the transformation order by order in $\gs$, introducing one auxiliary field at a time.
The purpose is to describe the different ingredients which are needed for rewriting the action at all orders (\cref{p:sec:all-orders}).

\Cref{p:sec:hs:action} recalls the action of closed string field theory and fixes the notations.
The quartic order (\cref{p:sec:hs:quartic}) introduces $\Sigma_2$ and is worked out in full.
At the quintic order (\cref{p:sec:hs:quintic}), $\Sigma_2$ alone is still sufficient, but its saddle-point solution and its gauge transformation receive corrections.
At the sextic order (\cref{p:sec:hs:sextic}), a $3$-string field $\Sigma_3$ is introduced, together with the reduced sextic vertex: the sextic vertex minus the exchange diagram induced by $\Sigma_2$.
The septic and octic orders (\cref{p:sec:hs:septic,p:sec:hs:octic}) make the general pattern visible.

We treat the quartic order in more detail than its simplicity requires, because every ingredient of the construction already appears there.
We start by describing the equation of motion which identifies the saddle-point of the auxiliary field with the tensor power $\Psi^2$, and the kernel of the quartic vertex with the quotient on which the auxiliary field lives.
We then show in detail how to eliminate the auxiliary field by completing a square.
We describe different properties of the field: its gauge transformation, its reality condition and its spacetime content.
Last come the Feynman rules through which the exchange of the auxiliary field reproduces the quartic vertex.
Each of these points returns, with corrections in $\gs$, at every higher order.

